% Nouveau EGA V (cote 25): archive cote 25, site page 122 (PDF page 123).
% typescript. Draft compared: none.
% AI-transcribed by Claude Opus 5.5 (claude-opus-5-5, adaptive thinking, effort high), Message Batches API, batch ega25b (msgbatch_01A7n5YuXubUPNK1zsepi8RG), 2026-11-13; not yet checked by eye.
\scanpage{25/122}
\mspage{25}

systêmes linéaires de diviseurs sans composante fixe  sur X paramétrés
par $\overset{\hand{\vee}}{P}$ , et l'ensemble Ø  des applications rationelles de X dans P, \hand{qui}
\var{\inserted{sur $\bar{k}$ ne se factorisent par aucun hyperplan de $P_{\unsure{\bar{k}}}$.}}{sur $\bar{k}$ ne se factorisent par aucun hyperplan de $P_{\unsure{\bar{k}}}$.}{insertion}{ega25_25_122-v0001}
\var{\inserted[typed]{lorsque les anneaux locaux de X sont factoriels.}}{lorsque les anneaux locaux de X sont factoriels.}{insertion}{ega25_25_122-v0002}

En effet, \var{\deleted{\var{\xout{X étant régulier, X en est de même de \unsure{X $X_S$} qui est lisse}}{}{typewriter strike-out}{ega25_25_122-v0004}
\var{\xout{sur X, donc d'après un résultat de Auslander-Buchsbaum (qui sera repris}}{}{typewriter strike-out}{ega25_25_122-v0005}
\var{\xout{au Chap III) \illegible{1} de les anneaux locaux de X\unsure{x}P sont factoriels,}}{}{typewriter strike-out}{ega25_25_122-v0006}
et comme codim(X-U,X) \quad 2}}{}{deletion}{ega25_25_122-v0003} , il s'ensuit que tout Module inversible sur
U peut se prolonger en un Module inversible sur X, donc la condition
envisagée dans b) est automatiquement satisfaite. D'autre part, d'après
\var{\deleted{\unsure{un Mémoire inédit lorsque n° ultérieur}}}{}{deletion}{ega25_25_122-v0007}
Auslander-Buchsbaum (\var{\deleted{qui figurera un jour au Chap III}}{}{deletion}{ega25_25_122-v0008}\var{\inserted{cf §20}}{cf §20}{insertion}{ega25_25_122-v0009}) un annea\hand{/}u local
régulier est factoriel.

\var{\deleted{La proposition 2  b)  nous montre que myennant une hypothèse}}{}{deletion}{ega25_25_122-v0010}
\var{\deleted{\var{\inserted[typed]{(savoir D \quad Ø)}}{(savoir D \quad Ø)}{insertion}{ega25_25_122-v0012}}}{}{deletion}{ega25_25_122-v0011}
\var{\deleted{supplémentaire sur un\var{\xout{e famille}}{}{typewriter strike-out}{ega25_25_122-v0014} système linéaire de diviseurs D, on peut}}{}{deletion}{ega25_25_122-v0013}
\var{\deleted{lui associer canoniquement un Module inversible $\underline{L}$ sur X, prolongement}}{}{deletion}{ega25_25_122-v0015}
\var{\deleted{\var{\inserted[typed]{où U=X-Z.}}{où U=X-Z.}{insertion}{ega25_25_122-v0017}}}{}{deletion}{ega25_25_122-v0016}
\var{\deleted{canonique de $\underline{L}^{x}_{U} = f\ (\underline{O}_P(1))$ ,  On peut en donner une autre interprétation,}}{}{deletion}{ega25_25_122-v0018}
\var{\deleted{indépendante de l'hypothèse  D \quad Ø , de la façon suivante; par constructio}}{}{deletion}{ega25_25_122-v0019}
\var{\deleted{\var{\inserted[typed]{can isomorphe à}}{can isomorphe à}{insertion}{ega25_25_122-v0021} \hand{)}}}{}{deletion}{ega25_25_122-v0020}
\var{\deleted{\var{\xout{l'image inverse de $\underline{L}$x}}{}{typewriter strike-out}{ega25_25_122-v0023} ou presque $\underline{L}^{x}_{U}\otimes_S\underline{O}_P(1)$ \var{\xout{n'}}{}{typewriter strike-out}{ega25_25_122-v0024}est \var{\xout{autre que}}{}{typewriter strike-out}{ega25_25_122-v0025} la restric-}}{}{deletion}{ega25_25_122-v0022}%
\var{\deleted{tion $\underline{N}'$ à $U\mathrm{x}_S\overset{\hand{\vee}}{P}$  du  Module inversible $\underline{N}$  sur $X\mathrm{x}_S\overset{\hand{\vee}}{P}$  défini par D. En vertu}}{}{deletion}{ega25_25_122-v0026}
\var{\deleted{du \var{\xout{\unsure{n°1}}}{}{typewriter strike-out}{ega25_25_122-v0028} n°1, cela détermine \var{\xout{uniquement $\underline{N}$, et}}{}{typewriter strike-out}{ega25_25_122-v0029} $\underline{L}_U$  à isom unique près;}}{}{deletion}{ega25_25_122-v0027}
\var{\deleted{et  (comme U est dense dans X)  il existe  un Module inversible $\underline{L}$ sur X,}}{}{deletion}{ega25_25_122-v0030}
\var{\deleted{déterminé à isom unique près (do\unsure{ç}c \unsure{prolongeant|prolongeaant} $\underline{L}_U$ ) tel que  l'on ait}}{}{deletion}{ega25_25_122-v0031}
\var{\deleted{un isomorphisme  $\underline{L}\otimes_{\underline{O}_S}\underline{O}_{\unsure{P}}\ (1) \quad \underline{N}$ ;  $\underline{L}$  est donné par exemple par}}{}{deletion}{ega25_25_122-v0032}
\var{\deleted{\[ \underline{L} \quad pr_1\ (\underline{N}) \quad ; \]}}{}{deletion}{ega25_25_122-v0033}
\var{\deleted{on l'appelle  le Module inversible sur X associé au systême lin de divi-}}{}{deletion}{ega25_25_122-v0034}%
\var{\deleted{suurs D.}}{}{deletion}{ega25_25_122-v0035}

\mspage{122}
